\section{本章小结}
\label{sec:density-summary}

本章从两个测度的比较出发，研究能否用一个逐点权重把参考测度换算成目标测度。第~\ref{chap:measures}章的绝对连续只说明零测集的包含关系；Radon–Nikodym定理表明，在两个测度都$\sigma$有限时，这个关系足以产生一个函数表示$\nu(A)=\int_Ah\,\mathrm d\mu$，并且该函数在$\mu$-几乎处处意义下唯一。证明的核心是在所有“不超过$\nu$”的候选权重中取积分最大者，再用Hahn分解说明剩余质量必为零。反过来，任何非负权重都能产生一个绝对连续的测度，这一方向不需要有限性条件。换测度公式把密度的作用从集合推广到函数，使目标测度下的积分化为参考测度下的加权积分。Lebesgue分解则保留了参考测度看不见的奇异质量：密度只表示绝对连续部分，点质量和Cantor测度这样的奇异连续质量都在密度之外。

使用密度时需要随时说明三件事。第一是参考测度：同一个目标测度相对于计数测度、几何权重或长度测度有不同的密度，不同参考下的密度不能直接相加，混合质量可以通过选择共同参考统一表示。第二是成立条件：存在性与唯一性需要$\sigma$有限，计数测度与长度的反例说明这一条件不能随意去掉；有符号测度要求总变差绝对连续，对应密度绝对可积；倒数公式只在测度等价时成立。第三是几乎处处：密度在参考零测集上的取值没有意义，链式规则与倒数公式都只在相应测度的几乎处处意义下成立。

离散质量、函数增量和映射后的质量是计算密度的三种常见入口。离散情形逐点相除；增量测度相对于长度的密度是导数，函数的绝对连续与测度的绝对连续在这里合而为一；推前测度的密度由Jacobian因子给出，多个分支的贡献需要累加。接下来的概率空间部分将在这些基础上继续：把一个随机量相对于较粗的信息作平均，需要在一个子$\sigma$代数上比较两种测度并取密度，条件期望正是这样定义的；比较两个概率分布的各种差异，也从它们相对于共同参考的密度出发。
